
%\needspace{3\baselineskip}
\color{uniblue}{\subsection{Контроль цепей напряжения стороны НН трансформатора (КЦН НН)}}
\color{black}

\begin{enumerate}[label=\arabic{section}.\arabic{subsection}.\arabic*, labelsep=4pt, leftmargin=0pt, itemindent=57pt]

% ЗДЕСЬ ВЕТВЛЕНИЕ, если исполнение Т2 то вставляем первый абзац
\ifboolexpr{bool {isT2}}
{
    \item
    В составе устройства предусматриваются две функции <<КЦН НН1>> и <<КЦН НН2>> (для стороны НН1 и НН2 соответственно), ниже приводится общее описание работы функции <<КЦН НН>>. 
}
{}

\item
Функциональный блок <<КЦН НН>> служит для сигнализации неисправности цепей напряжения стороны НН трансформатора.
\item
Ввод функции в работу на этапе параметрирования устройства осуществляется программным переключателем SGF1 <<Ввод\_функции>>. Информация о введенной функции отображается выходным сигналом <<КЦН НН: Ввод>>.
\item
Функция <<КЦН НН>> может быть выведена из работы оперативно путем активации сигнала <<ОВ КЦН НН>>, что характеризуется активным состоянием сигнала <<КЦН НН: Оперативный вывод>> на выходе функции.
\item
Функция <<КЦН НН>> контролирует снижение максимального из линейных напряжений (\textsl{UAB, UBC, UCA}) ниже заданной уставки и повышение напряжения обратной последовательности (\textsl{U2}) выше заданной уставки, в случае выполнения любого из условий через выдержку времени таймера Т1 (уставка <<Тср>>) выдается сигнал неисправности ЦН.
\item
Алгоритм работы функции <<КЦН НН>> выполнен в соответствии с рисунком \ref{kznnn:img1}. В таблице \ref{kznnn:tbl1} приведены параметры, необходимые для настройки функции <<КЦН НН>>.

\vspace{3mm}
\begin{figure}[H]
\centering
\includegraphics[width=0.9\textwidth,height=0.9\textheight,keepaspectratio]{\fbpath/57. КЦН НН общ/_latex/img/kznnn.pdf}
\captionof{figure}{Функциональная схема <<КЦН НН>>}
\label{kznnn:img1}
\end{figure}


\end{enumerate}
